(1) Fresnel传播核为 $e^{ikz}e^{ikr^2/(2z)}/(i\lambda z)$，圆盘外透光区域 $r>a$，故 $U_{\rm disk}(0,z)=E_0e^{ikz}(2\pi/(i\lambda z))\int_a^\infty r e^{ikr^2/(2z)}dr$。积分应加微弱正收敛因子后取极限。
(2) 令 $\phi=ka^2/(2z)$。利用 $\int_a^\infty r e^{ikr^2/(2z)-\epsilon r^2}dr$ 的收敛极限，并以无障碍平面波 $U_{\rm free}=E_0e^{ikz}$ 归一，得 $U_{\rm disk}=U_{\rm free}e^{i\phi}$。因此轴心 $|U_{\rm disk}|^2=|U_{\rm free}|^2$，虽有圆盘遮挡仍出现Poisson亮斑。
(3) Babinet互补性 $U_{\rm disk}+U_{\rm aperture}=U_{\rm free}$，故圆孔振幅 $U_{\rm aperture}=U_{\rm free}(1-e^{i\phi})$，强度 $I_{\rm aperture}=4I_0\sin^2(\phi/2)$。振幅相加而非强度相加，故两互补物体的强度一般不互补。
(4) 理想轴心亮斑要求单色、空间相干和圆对称边缘；扩展光源不同入射方向的亮斑错位并相互平均。还需 $a/z\ll1$ 的傍轴条件，有限孔径与探测器像素会改变可观测亮斑对比度。 当障碍尺寸与波长同阶时，还必须考虑矢量衍射和实际边缘形状。
